%! TEX root = part3.tex
% the main TeX file which is intended to compile, :VimtexReload after adjustment
% :h vimtex-tex-root
\documentclass[../main.tex]{subfiles}
% \includeonlyframes{current}
% \setbeameroption{hide notes} % Only slides
% \setbeameroption{show only notes} % Only notes
\setbeameroption{show notes on second screen=right} % Both

\begin{document}

\begin{frame}[mytitle=center,t,mycolor=digiPH_gray,light]
	\frametitle{Tugas Kelas (Lanjutan)}
	\txt{digiPH_leaf}{digiPH_white}{Buatlah maket Sloof berdasarkan denah yang tersedia.}

	\footnotesize Kelompok

\end{frame}


\begin{frame}[mytitle=center,t,mycolor=digiPH_gray,light]
	\frametitle{Tugas Rumah (Lanjutan)}
	\txt{digiPH_leaf}{digiPH_white}{Buatlah maket Kolom berdasarkan denah yang tersedia.}

	\footnotesize Kelompok

	Buatlah Gambar Kerja \Huge Rencana Pondasi \normalsize

	\footnotesize Individu

\end{frame}

\begin{frame}[mytitle=center,t,mycolor=digiPH_gray,light]
	\frametitle{Kolom}

	\begin{itemize}
		\item Kolom adalah elemen struktural vertikal yang berfungsi mentransfer beban dari struktur atas ke pondasi.
		\item Beban dari balok, lantai, dan atap diteruskan melalui kolom ke pondasi.
		\item Kolom berperan penting menjaga stabilitas bangunan terhadap:
		      \begin{itemize}
			      \item Beban vertikal
			      \item Beban lateral (gempa dan angin)
		      \end{itemize}
	\end{itemize}

\end{frame}


\begin{frame}[mytitle=standard,t,mycolor=digiPH_ubcblue,dark]
	\frametitle{Fungsi kolom:}
	\begin{itemize}
		\item Menyangga beban vertikal
		\item Menjaga kestabilan bangunan
		\item Menahan gaya horizontal
		\item Elemen arsitektural (estetika bangunan)
	\end{itemize}
\end{frame}



\begin{frame}[t,mycolor=digiPH_leaf,mytitle=standard,dark]
	\frametitle{Klasifikasi Kolom}
	\ltxt{digiPH_gray}{digiPH_darkblue}{\Huge Kolom \normalsize dapat diklasifikasikan berdasarkan:
		\begin{itemize}
			\item Material
			\item Bentuk penampang
			\item Jenis beban
			\item Posisi kolom
			\item Metode konstruksi
		\end{itemize}}
	\rimg{klasifikasi_kolom}
\end{frame}



\begin{frame}[label=current,t,mycolor=digiPH_leaf,mytitle=standard,dark]
	\frametitle{Penentuan Besaran Kolom}

	\begin{itemize}

		\item Penentuan ukuran kolom dapat menggunakan pendekatan bentangan modul.

		\item Modul adalah sistem grid yang digunakan untuk menentukan jarak antar kolom.

	\end{itemize}



\end{frame}




\begin{frame}[mytitle=center,c,mycolor=digiPH_ocean,dark]\frametitle{Perkiraan dimensi kolom:}
	\begin{itemize}
		\item $\frac{1}{10}$ sampai $\frac{1}{12}$ dari bentangan modul
		\item Bangunan 2 lantai dapat menggunakan sekitar $\frac{1}{20}$ dari bentangan modul
	\end{itemize}

	\Huge Contoh perhitungan: \normalsize

	\begin{itemize}
		\item Bentangan modul = 600 cm
		\item Dimensi kolom = $(1/10 \times 600)$
		\item Dimensi kolom = $60 \times 60$ cm
	\end{itemize}\end{frame}



\begin{frame}[t,mycolor=digiPH_leaf,mytitle=standard,dark]
	\frametitle{Komponen Kolom}


	Kolom umumnya terdiri dari beberapa \Huge elemen utama:
	\normalsize
	\begin{itemize}
		\item Penampang kolom (bentuk kolom)
		\item Tulangan baja (reinforcement)
		\item Bahan pengisi
		\item Penempatan kolom dalam struktur
	\end{itemize}
\end{frame}



\begin{frame}[t,mycolor=digiPH_leaf,mytitle=standard,dark]
	\frametitle{Parameter Perancangan Kolom}

	\framesubtitle{Beban yang ditanggung:}

	\begin{itemize}
		\item Beban mati
		\item Beban hidup
		\item Beban lateral (angin, gempa, tekanan tanah)
	\end{itemize}



\end{frame}


\begin{frame}[mytitle=standard,c,mycolor=digiPH_gray,light]\framesubtitle{Dimensi kolom:}

	\begin{itemize}
		\item Tinggi kolom (kekakuan dan risiko tekuk)
		\item Penampang kolom (daya dukung struktur)
	\end{itemize}
\end{frame}

\begin{frame}[mytitle=center,c,mycolor=digiPH_ocean,dark]
	\framesubtitle{Material struktur:}

	\begin{itemize}
		\item Kuat tekan beton ($f'_c$)
		\item Kuat leleh baja tulangan ($f_y$)
	\end{itemize}
\end{frame}



\begin{frame}[t,mycolor=digiPH_leaf,mytitle=standard,dark]
	\frametitle{Perbandingan Ukuran Kolom}

	\begin{table}
		\centering
		\begin{tabular}{c c c c}
			\textbf{No}  &  \textbf{Dimensi}  &  \textbf{Kuat Tekan}  &  \textbf{Tulangan}  \\
			             &  (cm)              &  (kN)                 &  \\
			\hline
			1            &  20 x 20           &  400                  &  4D12  \\
			2            &  25 x 25           &  625                  &  4D16  \\
			3            &  30 x 30           &  900                  &  6D16  \\
			4            &  40 x 40           &  1600                 &  8D20  \\
		\end{tabular}
	\end{table}

\end{frame}



\begin{frame}[t,mycolor=digiPH_leaf,mytitle=standard,dark]
	\frametitle{Perhitungan Luas Penampang Kolom}

	\begin{equation}
		A_c = \frac{P_u}{\phi f'_c}
	\end{equation}

	\begin{itemize}

		\item $A_c$ = luas penampang beton tekan
		\item $P_u$ = beban terfaktor
		\item $\phi$ = faktor reduksi kekuatan
		\item $f'_c$ = kuat tekan beton

	\end{itemize}

\end{frame}


\begin{comment}

\begin{frame}[t,mycolor=digiPH_leaf,mytitle=standard,dark]
	\frametitle{Perhitungan Tulangan Kolom}

	\textbf{Tulangan longitudinal}

	\begin{equation}
		A_s = \frac{P_u - \phi A_c f'_c}{\phi f_y}
	\end{equation}

	\begin{itemize}
		\item $A_s$ = luas tulangan baja
		\item $f_y$ = kuat leleh baja
	\end{itemize}

	\vspace{0.4cm}

	\textbf{Tulangan sengkang}

	\begin{equation}
		S = \frac{0.75 d}{\sqrt{f'_c}}
	\end{equation}

	\begin{itemize}
		\item $S$ = jarak antar sengkang
		\item $d$ = tinggi efektif penampang
	\end{itemize}

\end{frame}



\begin{frame}[t,mycolor=digiPH_leaf,mytitle=standard,dark]
	\frametitle{Pemeriksaan Stabilitas Kolom}

	\begin{itemize}

		\item Stabilitas kolom diperiksa terhadap kemungkinan tekuk.

	\end{itemize}

	\begin{equation}
		k = \frac{L_e}{r}
	\end{equation}

	\begin{itemize}

		\item $k$ = rasio kelangsingan kolom
		\item $L_e$ = panjang efektif kolom
		\item $r$ = jari-jari girasi penampang

		\item Rasio kelangsingan menentukan apakah kolom termasuk kolom pendek atau kolom langsing.

	\end{itemize}

\end{frame}


\end{comment}

\end{document}
